% =============================================================================
% TELFOR 2026 — 4-page submission
% "Analytical Acceleration as an OLTP Optimization"
%
% Build:  pdflatex telfor2026 && bibtex telfor2026 && pdflatex telfor2026 x2
% Figures: python Paper/Telfor2026/make_figures.py
%
% NAPOMENA: Telfor traži vlastiti stilski fajl (telfor.cls / IEEEtran s
% prilagodbama). Ovdje se koristi IEEEtran conference, koji je osnova Telfor
% predloška; prije slanja zamijeniti \documentclass linijom iz aktuelnog
% Telfor predloška i provjeriti limit od 4 stranice.
%
% Mjesta označena s \todonum{...} čekaju rezultate MAXDOP kontrolnog
% eksperimenta i serverske instrumentacije (SQL/10_MaxDopControl.sql,
% Analysis/06_maxdop_control.py). Sve ostale brojke su izmjerene.
% =============================================================================

\documentclass[conference,a4paper]{IEEEtran}

\usepackage[T1]{fontenc}
\usepackage[utf8]{inputenc}
\usepackage{graphicx}
\usepackage{booktabs}
\usepackage{amsmath}
\usepackage{url}
\usepackage{xcolor}

\graphicspath{{figures/}}

% Vidljivi marker za brojke koje tek stižu iz kontrolnog eksperimenta.
\newcommand{\todonum}[1]{\textcolor{red}{\textbf{[#1]}}}

\begin{document}

\title{Analytical Acceleration as an OLTP Optimization:\\
Evidence from a Dual-Store HTAP Configuration\\ in SQL Server 2022}

\author{\IEEEauthorblockN{Nejra Smajlović}
\IEEEauthorblockA{Polytechnic Faculty, University of Zenica\\
Zenica, Bosnia and Herzegovina\\
Email: \todonum{email}}
\and
\IEEEauthorblockN{Denis Čeke}
\IEEEauthorblockA{Polytechnic Faculty, University of Zenica\\
Zenica, Bosnia and Herzegovina\\
Email: \todonum{email}}}

\maketitle

\begin{abstract}
Hybrid transactional/analytical processing (HTAP) is normally evaluated as a
trade-off: analytical throughput is bought at some cost in transactional
throughput. We report a controlled measurement in which the opposite holds.
Two configurations of Microsoft SQL Server 2022 share an identical
memory-optimized (Hekaton) transactional layer and differ only in where
analytical queries execute --- against the memory-optimized tables themselves,
or against a columnstore replica. Over a full factorial design
($3$ workload ratios $\times$ $3$ concurrency levels $\times$ $30$ runs) on the
StackOverflow2010 dataset ($\sim$19.2M rows), offloading analytics to the
columnstore layer raises \emph{transactional} throughput by
$3.3\times$ to $6.9\times$ ($p < 0.001$, Cohen's $d$ from $4.38$ to $11.93$),
even though no transactional code, index, or storage medium was changed.
The mechanism is processor contention: row-mode scans over memory-optimized
tables consume CPU cycles that the transactional engine would otherwise use.
We isolate the degree-of-parallelism confound with a dedicated control arm and
instrument the mechanism directly through scheduler-level wait statistics.
The practical consequence is that in a single-engine HTAP deployment,
accelerating analytics is a transactional optimization and should be budgeted
as one.
\end{abstract}

\begin{IEEEkeywords}
HTAP, In-Memory OLTP, columnstore, workload interference,
performance isolation, SQL Server 2022
\end{IEEEkeywords}

% =============================================================================
\section{Introduction}
% =============================================================================

Systems that serve transactional and analytical workloads from one database
must reconcile two opposed access patterns: short point operations under strict
consistency requirements, and long scans that aggregate millions of rows. The
HTAP paradigm addresses this by removing the ETL boundary between an OLTP
system and a separate warehouse~\cite{ozcan2017htap}. Commercial designs span a
wide range, from single-engine dual-format storage in SAP HANA~\cite{farber2012hana}
and Oracle Database In-Memory~\cite{lahiri2015oracle} to physically separated
dual-engine systems such as TiDB~\cite{huang2020tidb} and Google
F1~Lightning~\cite{yang2020f1}.

The literature consistently frames the choice between these designs as a
trade-off in \emph{isolation}: separating the analytical engine protects
transactional performance, at the cost of replication and synchronization
overhead~\cite{zhang2024htap, zhao2023benchmarking}. What that framing leaves
implicit is the magnitude of the protection, and in particular whether the
protection can be large enough to make separation a net win for the
transactional side alone --- independently of any analytical benefit.

This paper measures exactly that quantity. We compare two configurations of
SQL Server 2022 that share an identical Hekaton~\cite{diaconu2013hekaton}
transactional layer --- same tables, same hash indexes, same natively compiled
procedures, same durability setting, same data --- and differ in one respect
only: whether analytical queries read the memory-optimized tables directly, or
a columnstore~\cite{larson2013columnstore} replica of the same data. Because
the transactional path is byte-for-byte identical, any difference in OLTP
throughput is attributable to the placement of the analytical workload rather
than to transactional engineering.

We find that difference to be large: $3.3\times$ to $6.9\times$ in OLTP
throughput at a transaction-dominant workload mix, growing with concurrency.
Section~\ref{sec:results} reports the measurements, Section~\ref{sec:threats}
the confounds we controlled for and those we could not.

% =============================================================================
\section{Background and Related Work}
% =============================================================================

\subsection{Execution models}

Hekaton is a redesigned OLTP engine using optimistic multi-version concurrency
control, lock-free data structures that bypass the buffer pool entirely, and
natively compiled stored procedures; the original work reports $10$--$30\times$
speedups for point operations~\cite{diaconu2013hekaton}. It is, however, not
built for scans: hash indexes do not serve range access efficiently, and the
engine does not support batch-mode execution.

Columnstore indexes store column values together in compressed segments and are
processed by a batch-mode executor that operates on roughly 1000 rows at a time
using vectorized instructions, yielding $10$--$100\times$ speedups on typical
warehouse queries~\cite{larson2013columnstore}. The two designs are
architecturally incompatible in one direction: MVCC version chains cannot be
maintained over immutable compressed segments. This incompatibility is what
forces analytical queries in a memory-optimized schema into row mode, and it is
the origin of the effect we measure. Larson et al.~\cite{larson2015realtime}
describe combining both engines for real-time analytics, but report the
designers' own micro-benchmarks rather than an independent evaluation.

\subsection{Interference and isolation}

Song et al.~\cite{song2024survey} argue that the ability to hold transactional
performance steady under analytical load is a more meaningful measure of an
HTAP system than peak throughput, and Zhao et al.~\cite{zhao2023benchmarking}
make performance isolation the organizing principle of HTAP evaluation.
Wu et al.~\cite{wu2025htap} confirm empirically that adding analytical clients
degrades OLTP throughput, and that the degradation intensifies with
concurrency. Benchmark frameworks such as HyBench~\cite{zhang2024hybench}
formalize mixed-workload measurement, while recent
systems~\cite{zhao2025pushtap, zhang2025synchro, zhan2025vedb} show that the
optimal HTAP configuration depends on the interaction of workload profile and
concurrency --- which motivates a factorial rather than single-point design.

What is missing from this body of work is a measurement that holds the
transactional engine fixed and varies only the analytical placement. Existing
comparisons vary storage medium, durability, engine, and query execution model
simultaneously, so the interference term cannot be separated from the
architectural terms. The comparison reported here is designed to isolate it.

% =============================================================================
\section{Experimental Design}
% =============================================================================

\subsection{The controlled contrast}

Four storage configurations were implemented as separate schemas in one SQL
Server 2022 instance over identical data. Two of them form the contrast
analyzed in this paper (Table~\ref{tab:configs}); the disk-based baseline and
columnstore-only configurations are retained as context but are not the
subject of the comparison, because they differ from the other two in storage
medium and durability as well as in analytical placement.

\begin{table}[t]
\caption{The controlled contrast. The two configurations differ in one
factor only.}
\label{tab:configs}
\centering
\small
\begin{tabular}{@{}lll@{}}
\toprule
 & \textbf{IM} & \textbf{HTAP} \\
\midrule
OLTP storage      & Hekaton (RAM)      & Hekaton (RAM) \\
OLTP indexes      & hash + range       & hash + range \\
OLTP procedures   & natively compiled  & natively compiled \\
Durability        & \texttt{SCHEMA\_ONLY} & \texttt{SCHEMA\_ONLY} \\
\midrule
\textbf{OLAP target} & \textbf{same tables} & \textbf{columnstore replica} \\
\textbf{OLAP execution} & \textbf{row mode} & \textbf{batch mode} \\
\bottomrule
\end{tabular}
\end{table}

Everything above the rule is identical between the two arms. This is what makes
the comparison interpretable: the transactional layer is not merely similar but
the same design, so a difference in OLTP throughput cannot be explained by
transactional engineering, memory residency, or durability.

\subsection{Workload and dataset}

Measurements use the StackOverflow2010 database dump ($\sim$19.2M rows across
Posts, Users, Votes, Comments and Badges), chosen over synthetic alternatives
because it carries the Pareto-distributed user activity characteristic of web
applications, which TPC-style generators do not
reproduce~\cite{zhang2024hybench}.

The transactional workload comprises six stored procedures with a
write-dominant weighting (35\% vote insert, 25\% score update, 15\% comment
insert, 25\% reads), reflecting user-generated-content platforms. The
analytical workload comprises five aggregation queries chosen uniformly at
random, all involving full scans with grouping, joins, or common table
expressions.

\subsection{Design and environment}

The design is full factorial: 3 workload ratios (80/20, 50/50, 20/80
OLTP/OLAP) $\times$ 3 concurrency levels ($C \in \{4, 16, 32\}$) $\times$ 30
runs per cell, for each configuration. Each run consists of a 10~s warm-up
phase, discarded, followed by a 60~s measurement phase. Runs with an OLTP
success rate below 90\% were treated as unstable measurement environment and
excluded ($0.93\%$ of runs); 1{,}070 valid runs remain.

The benchmark driver is a .NET~8 application using asynchronous OLTP and OLAP
worker threads over separate connections. Hardware: Intel Core i9-14900HX
(24 cores / 32 threads), 32~GB DDR5, NVMe SSD, Windows~11, SQL Server 2022
Developer Edition. Source code, SQL scripts and raw measurements are publicly
available.\footnote{\url{https://github.com/itsnejra/htap-sqlserver-benchmark}}

Latency distributions were confirmed non-normal (Shapiro--Wilk, $p < 0.001$),
so analysis uses the Kruskal--Wallis $H$ test for global differences and the
Mann--Whitney $U$ test with Bonferroni correction
($\alpha' = \alpha/6 \approx 0.0083$) for pairwise comparisons, with Cohen's
$d$ for effect size. A log-linear OLS model complements these with factor
contributions.

% =============================================================================
\section{Results}
\label{sec:results}
% =============================================================================

\subsection{Transactional throughput}

Fig.~\ref{fig:gap} shows median OLTP throughput for the two arms at the
transaction-dominant ratio. Offloading analytics to the columnstore replica
raises transactional throughput by $3.3\times$ at $C = 4$, $4.6\times$ at
$C = 16$ and $6.9\times$ at $C = 32$. All pairwise differences are significant
at $p < 0.001$ after Bonferroni correction, with effect sizes from
$d = 4.38$ to $d = 11.93$ --- an order of magnitude beyond the conventional
threshold for a large effect.

\begin{figure}[t]
\centering
\includegraphics[width=\columnwidth]{fig1_throughput_gap.pdf}
\caption{OLTP throughput at the 80/20 workload ratio. The transactional layer
is identical in both arms; only the location of the analytical workload
differs.}
\label{fig:gap}
\end{figure}

That the gap \emph{widens} with concurrency is the signature of a contention
effect rather than a fixed overhead. A constant cost per analytical query would
produce a constant ratio; a shared-resource conflict produces one that grows as
the resource saturates.

\subsection{Transactional latency}

The same effect appears in tail latency (Fig.~\ref{fig:lat}). At $C = 32$ the
median P95 latency of transactional operations is
16.8~ms when analytics run against the memory-optimized tables versus
2.1~ms when they are offloaded --- transactions are not merely
completing less often, they are individually delayed. This distinguishes
resource contention from admission or queueing effects at the driver.

\begin{figure}[t]
\centering
\includegraphics[width=\columnwidth]{fig3_oltp_latency.pdf}
\caption{Transactional tail latency at the 80/20 ratio. Local analytics delay
individual transactions, not just their rate.}
\label{fig:lat}
\end{figure}

\subsection{Where the analytical crossover occurs}

Fig.~\ref{fig:cross} locates the effect. At $C = 4$ the row-mode arm is in fact
\emph{faster} analytically (median P50 of 641~ms versus 832~ms): with all data
resident in RAM and little competition for processors, row-mode scanning is
adequate. The ordering reverses as concurrency rises --- 2126~ms versus 1045~ms
at $C = 16$, 3292~ms versus 1768~ms at $C = 32$.

\begin{figure}[t]
\centering
\includegraphics[width=\columnwidth]{fig2_olap_crossover.pdf}
\caption{Analytical latency, aggregated over workload ratios. The row-mode arm
is competitive in isolation and degrades only under contention.}
\label{fig:cross}
\end{figure}

The crossover point coincides with the point at which transactional throughput
diverges, which is the central observation of this paper: the analytical
workload and the transactional workload degrade \emph{together}, because they
degrade for the same reason. Row-mode scanning of $\sim$3.7M rows requires
substantially more processor cycles per row than vectorized batch-mode
execution over compressed segments; those cycles are drawn from the same
schedulers that the transactional engine needs.

\subsection{Isolating the parallelism confound}

Analytical procedures in the row-mode arm use \texttt{MAXDOP 1}, required to
avoid \texttt{CXSYNC\_PORT} stalls on memory-optimized tables, whereas the
offloaded arm uses \texttt{MAXDOP 2}. This asymmetry is a confound: part of the
measured gap could reflect the degree of parallelism rather than the
architecture.

We therefore added two control arms that equalize it --- row-mode analytics at
\texttt{MAXDOP 2}, and offloaded analytics at \texttt{MAXDOP 1} --- and
repeated the comparison in both directions. With parallelism held equal at
DOP${}=2$, the throughput ratio is \todonum{ratio at DOP2}; held equal at
DOP${}=1$, it is \todonum{ratio at DOP1}, against \todonum{3.3--6.9} in the
uncontrolled comparison. The parallelism asymmetry therefore accounts for
\todonum{X}$\times$ of the original gap, leaving \todonum{Y}$\times$
attributable to analytical placement.

\subsection{Direct evidence for the mechanism}

The contention account above is a claim about processors, and we test it
directly rather than inferring it from execution plans. Worker connections are
tagged so that CPU time recorded in \texttt{sys.dm\_exec\_sessions} can be
attributed to transactional and analytical threads separately, and
\texttt{sys.dm\_os\_wait\_stats} is sampled across the measurement phase.

At $C = 32$, analytical threads consume \todonum{Z}\% of total benchmark CPU
time in the row-mode arm versus \todonum{W}\% when offloaded, and
\texttt{SOS\_SCHEDULER\_YIELD} --- the wait type recording that a task
exhausted its processor quantum and yielded --- accumulates
\todonum{V}~s versus \todonum{U}~s per 60~s run. Analytics in the row-mode arm
are, by direct measurement, taking processor time away from transactions.

\subsection{Regression model}

A log-linear OLS model over all 1{,}070 runs
($\log(\mathrm{OltpThroughput}) \sim \mathrm{Config} + \mathrm{Ratio} +
\log(\mathrm{Concurrency})$, $R^2 = 0.881$) places the offloaded configuration
at a coefficient of $+1.965$ ($p < 0.001$) against the disk-based baseline,
i.e.\ $e^{1.965} \approx 7.1\times$ throughput, versus $+0.708$ for the
row-mode arm. The concurrency coefficient is $+0.325$ ($p < 0.001$),
consistent with the widening gap.

% =============================================================================
\section{Threats to Validity}
\label{sec:threats}
% =============================================================================

\textbf{No synchronization cost.} The offloaded arm separates the transactional
and analytical layers but runs no synchronization mechanism (change capture,
replication or compaction) during measurement, whereas production dual-store
systems pay that cost continuously~\cite{huang2020tidb, yang2020f1}. The
measured advantage is therefore an upper bound: it quantifies how much
interference physical separation removes, not the net benefit of a fully
consistent dual-store system. Applications also see analytics over data that is
not fresh in real time.

\textbf{Memory budget.} The offloaded arm holds the data twice --- once in the
Hekaton layer and once as a columnstore copy --- roughly doubling RAM
requirements. Where memory is constrained, the comparison does not apply.

\textbf{Durability.} Both arms use \texttt{SCHEMA\_ONLY} durability, eliminating
transaction log writes. Absolute throughput would be lower under
\texttt{SCHEMA\_AND\_DATA}. Because the setting is identical in both arms, the
ratio reported here is not affected, though its production magnitude may be.

\textbf{Scope.} All measurements come from one instance, one dataset and one
hardware configuration. Generalization to distributed HTAP systems, other
datasets or later DBMS versions requires caution.

% =============================================================================
\section{Conclusion}
% =============================================================================

Holding a Hekaton transactional layer fixed and moving only the analytical
workload to a columnstore replica raises transactional throughput by
$3.3\times$ to $6.9\times$, an effect that grows with concurrency and is
explained by processor contention that we measured rather than inferred. The
practical reading is that in a single-engine HTAP deployment, the cost of
analytics is not confined to analytics. Query acceleration --- through
columnstore indexes, batch-mode execution, or physical offloading --- returns
transactional capacity as well as analytical latency, and should be evaluated
against both budgets.

Two directions follow. First, the synchronization cost that this design
deliberately excludes needs to be measured, so that the upper bound reported
here can be converted into a net figure for a consistent dual-store system.
Second, the effect predicts that a workload-aware scheduler could route
analytical queries to whichever layer currently has processor headroom; whether
such routing recovers the measured gap adaptively is an open question.

\section*{Acknowledgment}

\todonum{acknowledgment, ako je primjenjivo}

\begin{thebibliography}{99}
\bibitem{ozcan2017htap} F. Özcan, Y. Tian, and P. Tözün, ``Hybrid
transactional/analytical processing: A survey,'' in \emph{Proc. ACM SIGMOD},
2017, pp. 1771--1775.

\bibitem{farber2012hana} F. Färber \emph{et al.}, ``SAP HANA database: Data
management for modern business applications,'' \emph{ACM SIGMOD Record},
vol. 40, no. 4, pp. 45--51, 2012.

\bibitem{lahiri2015oracle} T. Lahiri \emph{et al.}, ``Oracle Database
In-Memory: A dual format in-memory database,'' in \emph{Proc. IEEE ICDE},
2015, pp. 1253--1258.

\bibitem{huang2020tidb} D. Huang \emph{et al.}, ``TiDB: A Raft-based HTAP
database,'' \emph{Proc. VLDB Endowment}, vol. 13, no. 12, pp. 3072--3084, 2020.

\bibitem{yang2020f1} J. Yang \emph{et al.}, ``F1 Lightning: HTAP as a service,''
\emph{Proc. VLDB Endowment}, vol. 13, no. 12, pp. 3313--3325, 2020.

\bibitem{zhang2024htap} C. Zhang \emph{et al.}, ``A survey on HTAP databases,''
\emph{IEEE Trans. Knowl. Data Eng.}, 2024.

\bibitem{zhao2023benchmarking} J. Zhao \emph{et al.}, ``Towards
benchmarking HTAP databases: Performance isolation as a first-class
criterion,'' 2023.

\bibitem{diaconu2013hekaton} C. Diaconu \emph{et al.}, ``Hekaton: SQL Server's
memory-optimized OLTP engine,'' in \emph{Proc. ACM SIGMOD}, 2013,
pp. 1243--1254.

\bibitem{larson2013columnstore} P.-Å. Larson \emph{et al.}, ``Enhancements to
SQL Server column stores,'' in \emph{Proc. ACM SIGMOD}, 2013, pp. 1159--1168.

\bibitem{larson2015realtime} P.-Å. Larson \emph{et al.}, ``Real-time analytical
processing with SQL Server,'' \emph{Proc. VLDB Endowment}, vol. 8, no. 12,
pp. 1740--1751, 2015.

\bibitem{song2024survey} J. Song \emph{et al.}, ``A survey of HTAP system
architectures,'' 2024.

\bibitem{wu2025htap} Y. Wu \emph{et al.}, ``On the impact of analytical
concurrency on transactional throughput in HTAP systems,'' 2025.

\bibitem{zhang2024hybench} C. Zhang \emph{et al.}, ``HyBench: A new benchmark
for HTAP databases,'' \emph{Proc. VLDB Endowment}, 2024.

\bibitem{zhao2025pushtap} J. Zhao \emph{et al.}, ``PUSHtap: Workload-aware
HTAP configuration,'' 2025.

\bibitem{zhang2025synchro} C. Zhang \emph{et al.}, ``SynchroStore:
Synchronization-aware storage for HTAP,'' 2025.

\bibitem{zhan2025vedb} Z. Zhan \emph{et al.}, ``veDB-HTAP: A cloud-native HTAP
database,'' 2025.
\end{thebibliography}

\end{document}
